\documentstyle[%


righttag%


,amssymb%


,11pt%


,russian%


,izvru%


]{amsart}





%\setlength{\textwidth}{170mm}


%\setlength{\textheight}{232mm}


%\setlength{\hoffset}{-18mm}


%\setlength{\voffset}{-12mm}


%\setlength{\parindent}{6mm}





\newcommand{\tl}[3]{\noindent{\bf#1}\ #2 \dotfill #3 \par} 


\setcounter{page}{91}


\pagestyle{plain}


\begin{document}


\par


\vskip 2cm


\par


\bigskip


\bigskip


\bigskip


\par


{\centerline{\Large СОДЕРЖАНИЕ ЖУРНАЛА ``ИЗВЕСТИЯ ВЫСШИХ УЧЕБНЫХ ЗАВЕДЕНИЙ''}}





{\centerline{\Large серия ``Математика'' за 2002 г.}}


\bigskip


\bigskip








\tl{Абдурагимов Э.И.}


{О единственности положительного решения одной нелинейной\newline


двухточечной краевой задачи}{\No\,6}


\tl{Авсянкин О.Г.}


{О применении проекционного метода к парным интегральным операторам\newline


с однородными ядрами}{\No\,8}


\tl{Айзенберг Л.А., Гроссман И.Б., Коробейник Ю.Ф.}


{Некоторые замечания о радиусе\newline Бора для степенных рядов}{\No\,10}


\tl{Акишев Г.А.}


{О некоторых теоремах для системы Прайса}{\No\,1}


\tl{Аксентьев Л.А.} 


{Локальное строение поверхности внутреннего 


конформного радиуса для\newline плоской области}{\No\,4}


\tl{Андрианова А.А., Заботин Я.И.}


{Управление процессом минимизации \newline


в параметризованном методе центров}{\No\,12}


\tl{Антипина Н.В., Дыхта В.А.}


{Линейные функции Ляпунова--Кротова и достаточные\newline условия


оптимальности в форме принципа максимума}{\No\,12}


\tl{Аргучинцев А.В.}


{Решение задачи оптимального управления начально-краевыми\newline


условиями гиперболической системы на основе точных


формул приращения}{\No\,12}


\tl{Аргучинцева М.А.}


{Вариационная задача об оптимальной форме тела


с минимальным\newline радиационным нагревом поверхности}{\No\,2}


\tl{Ащепков Л.Т., Давыдов Д.В.}


{Стабилизация наблюдаемой линейной системы\newline управления


с постоянными интервальными коэффициентами}{\No\,2}


\tl{Бадам У.}


{Необходимые условия оптимальности в задачах живучести}{\No\,2}


\tl{Банару М.Б.}


{Две теоремы о косимплектических гиперповерхностях


6-мерных эрмитовых\newline подмногообразий алгебры Кэли}{\No\,1}


\tl{Бариева Н.А., Бикчантаев И.А.} 


{Задача сопряжения для неоднородного уравнения\newline


Коши--Римана на римановой поверхности}{\No\,8}


\tl{Бикчантаев И.А., Бариева Н.А.} 


{Задача сопряжения для неоднородного уравнения\newline


Коши--Римана на римановой поверхности}{\No\,8}


\tl{Билялов Р.Ф.}{Спиноры на римановых многообразиях}{\No\,11}


\tl{Бобочко В.Н.}


{Дифференциальная точка поворота


в теории сингулярных возмущений. I}{\No\,3}


\tl{Бобочко В.Н.}{Дифференциальная точка поворота


в теории сингулярных возмущений. II}{\No\,5}


\tl{Булдаев А.С.} 


{К оптимизации квадратичных по состоянию


динамических систем}{\No\,12}


\tl{\fbox{Васильев О.В.}, Лутковская Е.А.}


{Об одной 


процедуре улучшения управления\newline в параметрическом синтезе}{\No\,12}


\tl{Васильев С.Н.}


{Методы фрактальной интерполяции типа Барнсли}{\No\,9}


\tl{Вишневский B.В.}


{Вклад Бояи, Гаусса и Лобачевского в открытие неевклидовой


геометрии\newline (к 200-летию со дня рождения Яноша Бояи)}{\No\,11}


\tl{Власов В.В.} 


{О базисности семейства экспоненциальных  


решений дифференциально-разност-\newline ных уравнений}{\No\,6}


\tl{Bодопьянов C.K., Ухлов А.Д.}


{Операторы суперпозиции в пространствах Соболева}{\No\,10}


\tl{Волкодавов В.Ф., Илюшина Ю.А.} 


{Характеристический принцип локального\newline экстремума для одного 


уравнения гиперболического типа и его применение}{\No\,4}


\tl{Воронин А.Ф.}


{О корректности краевой задачи на прямой


для трех аналитических\newline функций}{\No\,4}


\tl{Воронин А.Ф.}


{Аналог теоремы Пикара для уравнения 


1-го рода в свертках с гладким\newline ядром}{\No\,7}


\pagebreak


\tl{Воронин С.М., Мещерякова Ю.И.}


{Аналитическая классификация типичных\newline вырожденных


элементарных особых точек ростков голоморфных векторных полей


на\newline комплексной плоскости}{\No\,1}


\tl{Габасов Р., Кириллова Ф.М.}


{Стабилизация систем с запаздываниями методами\newline


оптимального управления}{\No\,12}


\tl{Габдулхаев Б.Г.}


{Решение операторных и интегральных уравнений


методом\newline Боголюбова--Крылова}{\No\,10}


\tl{Габдулхаев Б.Г.}


{Оптимизация прямых методов решения 


периодических краевых задач}{\No\,12}


\tl{Габдулхаев Б.Г., Рахимов И.К.}


{Об одном оптимальном сплайн-методе решения\newline операторных


уравнений}{\No\,2}


\tl{Гайсин А.М., Латыпов И.Д.}


{Асимптотика логарифма максимального члена\newline измененного ряда Дирихле}{\No\,9}


\tl{Галиуллина Г.Р., Насыров С.Р.} 


{Уравнение Гахова для внешней смешанной 


обратной\newline краевой задачи по параметру $x$}{\No\,10}


\tl{Гарифьянов Ф.Н.} 


{Представляющие системы периодических 


мероморфных функций}{\No\,8}


\tl{Го Вэнь Бинь, Скиба А.Н.}{Два замечания о тождествах решеток


$\omega$-локальных и\newline $\omega$-композиционных формаций


конечных групп}{\No\,5}


\tl{Григорян М.Г.}


{Об одной ортонормированной системе}{\No\,4}


\tl{Гроссман И.Б., Айзенберг Л.А., Коробейник Ю.Ф.}


{Некоторые замечания о радиусе\newline Бора для степенных рядов}{\No\,10}


\tl{Гусаренко С.А.}


{Об условиях положительной определенности


операторов в гильбертовом\newline пространстве}{\No\,1}


\tl{Давыдов Д.В., Ащепков Л.Т.}


{Стабилизация наблюдаемой линейной системы\newline управления


с постоянными интервальными коэффициентами}{\No\,2}


\tl{Денисюк И.Т.} 


{Задача сопряжения гармонических функций 


в трехмерных областях\newline с негладкими границами}{\No\,4}


\tl{Дубровский В.В., Смирнова Л.В.} 


{К вопросу о восстановлении потенциала 


в обратной\newline задаче Робина}{\No\,7}


\tl{Дыхта В.А., Антипина Н.В.}


{Линейные функции Ляпунова--Кротова и достаточные\newline условия


оптимальности в форме принципа максимума}{\No\,12}


\tl{Ермолаев Ю.Б.}


{О пропорциональности лиевых слов 


в классических алгебрах Ли}{\No\,9}


\tl{Ермолаева Л.Б.}


{Решение интегральных уравнений


методом подобластей}{\No\,9}


\tl{Ермолицкий А.А.} 


{Теорема о шляпе и проблемы классификации 


структур на\newline римановых многообразиях}{\No\,11}


\tl{Жегалов В.И., Миронов А.Н.}


{О задачах Коши для двух уравнений в частных\newline производных}{\No\,5}


\tl{Железовский С.Е.}


{Оценки скорости сходимости


проекционно-разностного метода


для\newline гиперболических уравнений}{\No\,1}


\tl{Жогова Т.Б.} 


{О существовании семейств $R^m_{2n-1}$, допускающих проективное 


изгибание\newline второго порядка}{\No\,1}


\tl{Заботин Я.И., Андрианова А.А.}


{Управление процессом минимизации \newline


в параметризованном методе центров}{\No\,12}


\tl{Иванов А.Г.}


{О корректности расширения задач


управления почти периодическими\newline движениями}{\No\,6}


\tl{Игудесман К.Б.} 


{Фрактальная размерность пересечения стандартных канторовых \newline


множеств}{\No\,11}


\tl{Ильичев В.Г.} 


{Наследуемые свойства неавтономных динамических систем\newline


и их приложение в моделях конкуренции}{\No\,6}


\tl{Илюшина Ю.А., Волкодавов В.Ф.} 


{Характеристический принцип локального\newline экстремума для одного 


уравнения гиперболического типа и его применение}{\No\,4}


\tl{Исламов Г.Г.}


{О допустимых помехах линейных управляемых систем}{\No\,2}


\tl{Касапенко Л.Ю.}


{Некоторые алгоритмические вопросы ассоциативных алгебр}{\No\,11}


\tl{Кац Б.А.}


{Об обобщении одной теоремы Н.А.\,Давыдова}{\No\,1}


\tl{Кац Б.А., Погодина А.Ю.}


{О граничных значениях интеграла типа Коши на негладкой\newline


кривой}{\No\,3}


\tl{Кириллова Ф.М., Габасов Р.}


{Стабилизация систем с запаздываниями методами\newline


оптимального управления}{\No\,12}


\tl{Киясов С.Н.}{Некоторые классы сингулярных  


интегральных уравнений, разрешаемые\newline в замкнутой форме}{\No\,5}


\tl{Климов В.С.} 


{Симметризация функций из пространств Соболева}{\No\,11}


\tl{Коган Ю.В., Рубиновский А.В.}


{Исследование решений уравнения для дрейфов\newline


волнового твердотельного гироскопа


в зависимости от параметров}{\No\,6}


\tl{Кокурин М.Ю.}


{Условие истокопредставимости и оценки скорости 


сходимости методов\newline регуляризации линейных уравнений в 


банаховом пространстве. II}{\No\,3}


\tl{Коняев Ю.A.}


{Метод унитарных преобразований в теории устойчивости}{\No\,2}


\tl{Коняев Ю.А.} 


{Алгоритм построения квазирегулярного асимптотического 


представления\newline решения сингулярно возмущенных линейных многоточечных 


краевых задач с быстрыми\newline и медленными переменными}{\No\,7}


\tl{Корешков Н.А.} 


{Элементы Казимира $\Bbb Z$-форм модулярных алгебр Ли}{\No\,3}


\tl{Корешков Н.А.}


{Центральные элементы и инварианты в модулярных алгебрах Ли}{\No\,7}


\tl{Коробейник Ю.Ф., Айзенберг Л.А., Гроссман И.Б.}


{Некоторые замечания о радиусе\newline Бора для степенных рядов}{\No\,10}


\tl{Кузнецов А.А.}


{Представление однолистных функций  в виде бесконечных композиций}{\No\,10}


\tl{Лаптев Г.Г.}


{Об отсутствии решений


эллиптических дифференциальных неравенств\newline


в окрестности конической точки границы}{\No\,9}


\tl{Латыпов И.Д., Гайсин А.М.}


{Асимптотика логарифма максимального члена\newline измененного ряда Дирихле}{\No\,9}


\tl{Левенштам В.Б.} 


{Асимптотическое интегрирование квазилинейных параболических\newline


уравнений с быстро осциллирующими по времени коэффициентами}{\No\,5}


\tl{Ломп Х., Насрутдинов М.Ф., Сахаев И.И.}


{О проективных модулях с полулокальным\newline кольцом эндоморфизмов}{\No\,8}


\tl{Лутковская Е.А., \fbox{Васильев О.В.}}


{Об одной 


процедуре улучшения управления\newline в параметрическом синтезе}{\No\,12}


\tl{Маликов А.И.} 


{Эллипсоидальное оценивание решений дифференциальных уравнений\newline с 


помощью матричных систем сравнения}{\No\,8}


\tl{Малыгина В.В.} 


{Об оценке матрицы Коши некоторых систем с последействием}{\No\,6}


\tl{Марченко И.И.}


{Об оценках верхней и нижней логарифмических плотностей


в теории\newline мероморфных функций}{\No\,10}


\tl{Марюкова Н.Е.}


{Поверхности постоянной кривизны


в квазипсевдоримановом пространстве\newline


постоянной кривизны и уравнение Клейна--Гордона}{\No\,3}


\tl{Махер А., Плещинский Н.Б.}


{Задача о скачке для уравнения Гельмгольца\newline в плоскослоистой среде


и ее приложения}{\No\,1}


\tl{Мещерякова Ю.И., Воронин С.М.}


{Аналитическая классификация типичных\newline вырожденных


элементарных особых точек ростков голоморфных векторных полей


на\newline комплексной плоскости}{\No\,1}


\tl{Микенберг М.А.} 


{Некоторые топологические свойства 


многообразий над локальной\newline алгеброй, 


допускающих голоморфные вложения}{\No\,11}


\tl{Миронов А.Н., Жегалов В.И.}


{О задачах Коши для двух уравнений в частных\newline производных}{\No\,5}


\tl{Мирзаханян Э.А.} 


{Построение бесконечномерных относительных гомотопических групп\newline в 


гильбертовом пространстве}{\No\,8}


\tl{Мокейчев В.С.}


{Проблема собственных значений во всем пространстве


для уравнений\newline с разрывными коэффициентами}{\No\,6}


\tl{Муминов К.К.}


{Эквивалентность путей относительно действия


симплектической группы}{\No\,7}


\tl{Мухарлямов Р.Г.} 


{О численном решении уравнений экстремалей 


вариационной задачи\newline с ограничениями}{\No\,4}


\tl{Мысливец С.Г.}


{Граничное поведение интеграла типа логарифмического вычета}{\No\,4}


\tl{Насрутдинов М.Ф., Ломп Х., Сахаев И.И.}


{О проективных модулях с полулокальным\newline кольцом эндоморфизмов}{\No\,8}


\tl{Насыров С.Р., Галиуллина Г.Р.} 


{Уравнение Гахова для внешней смешанной 


обратной\newline краевой задачи по параметру $x$}{\No\,10}


\tl{Пинус А.Г.}{Об элементарной эквивалентности


производных структур свободных решеток}{\No\,5}


\tl{Плещинский Н.Б., Махер А.}


{Задача о скачке для уравнения Гельмгольца\newline в плоскослоистой среде


и ее приложения}{\No\,1}


\tl{Погодина А.Ю., Кац Б.А.}


{О граничных значениях интеграла типа Коши на негладкой\newline


кривой}{\No\,3}


\tl{Попов Ю.И.}


{Почти контактные структуры


многообразия $P^0_n(\cal H)$}{\No\,1}


\tl{Попова С.Н.}


{Об эквивалентности локальной достижимости


и полной управляемости\newline линейных систем}{\No\,6}


\tl{Ратыни А.К.} 


{Об эллиптической краевой задаче с оператором суперпозиции  


в граничном\newline условии. II}{\No\,6}


\tl{Рахимов И.К., Габдулхаев Б.Г.}


{Об одном оптимальном сплайн-методе решения\newline операторных


уравнений}{\No\,2}


\tl{Рубиновский А.В., Коган Ю.В.}


{Исследование решений уравнения для дрейфов\newline


волнового твердотельного гироскопа


в зависимости от параметров}{\No\,6}


\tl{Рузакова О.А., Федоров В.Е.}


{Управляемость линейных уравнений соболевского типа\newline


с относительно $p$-радиальными операторами}{\No\,7}


\tl{Сазанова Л.А.}


{Устойчивость оптимального синтеза


в задаче быстродействия}{\No\,2}


\tl{Салий B.H.}


{Идемпотентные полугруппы с транзитивным отношением


коммутативности}{\No\,1}


\tl{Сахаев И.И., Ломп Х., Насрутдинов М.Ф.}


{О проективных модулях с полулокальным\newline кольцом эндоморфизмов}{\No\,8}


\tl{Секерин А.Б.}


{Представление функций разностью


штейнеровских потенциалов}{\No\,8}


\tl{Семенов Е.М., Франкетти К.}


{Ортогональные подпространства перестановочно-инвари-\newline антных


пространств}{\No\,7}


\tl{Сихов М.Б.}


{Неравенства типа Бернштейна, Джексона--Никольского и 


их приложения}{\No\,8}


\tl{Скиба А.Н., Го Вэнь Бинь}{Два замечания о тождествах решеток


$\omega$-локальных и\newline $\omega$-композиционных формаций


конечных групп}{\No\,5}


\tl{Смирнова Л.В., Дубровский В.В.} 


{К вопросу о восстановлении потенциала 


в обратной\newline задаче Робина}{\No\,7}


\tl{Смольникова М.В.}{Обобщенно рекуррентное


симметрическое тензорное поле}{\No\,5}


\tl{Смольникова М.В., Степанов С.Е.}


{Фундаментальные дифференциальные операторы\newline


первого порядка на внешних и симметрических формах}{\No\,11}


\tl{Сосов Е.Н.}


{О непрерывности метрической $\delta$-проекции на 


выпуклое множество\newline в специальном метрическом


пространстве}{\No\,1}


\tl{Срочко В.А.}


{Модернизация методов градиентного типа 


в задачах оптимального\newline управления}{\No\,12}


\tl{Старков В.В., Якубовский З.Ю.}


{О совпадении двух линейно-инвариантных семейств\newline функций}{\No\,4}


\tl{Степанов С.Е., Смольникова М.В.}


{Фундаментальные дифференциальные операторы\newline


первого порядка на внешних и симметрических формах}{\No\,11}


\tl{Столяров А.В.}{Оснащения и аффинные связности 


на распределениях конформного\newline пространства}{\No\,5}


\tl{Столяров А.В.} 


{Внутренняя геометрия нормализованного конформного пространства}{\No\,11}


\tl{Тер\"ехин М.Т.}


{Ненулевые периодические решения неавтономной


системы обыкновенных\newline дифференциальных уравнений}{\No\,6}


\tl{Тронин С.Н.}


{Операды в категории конвексоров. I}{\No\,3}


\tl{Тронин С.Н.}{Операды в категории конвексоров. II}{\No\,5}


\tl{Угулава Д.К.}


{Об одном аналоге теоремы Пэли--Винера}{\No\,8}


\tl{Усач\"ев Ю.В.} 


{Бифуркация ограниченных решений системы обыкновенных 


дифференци-\newline альных уравнений}{\No\,7}


\tl{Ухлов А.Д., Bодопьянов C.K.}


{Операторы суперпозиции в пространствах Соболева}{\No\,10}


\tl{Федоров В.Е., Рузакова О.А.}


{Управляемость линейных уравнений соболевского типа\newline


с относительно $p$-радиальными операторами}{\No\,7}


\tl{Фернандеш В.У.}


{Новый класс делителей полугрупп изотонных отображений


конечных\newline цепей}{\No\,3}


\tl{Франкетти К., Семенов Е.М.}


{Ортогональные подпространства перестановочно-инвари-\newline антных


пространств}{\No\,7}


\tl{Хавинсон С.Я.} 


{Аппроксимация элементами клина с учетом  величин 


аппроксимирующих\newline элементов}{\No\,10}


\tl{Хайруллин Р.С.}


{Аналог задачи Франкля для уравнения второго рода}{\No\,4}


\tl{Чеголин А.П.}


{Классы $L_{p,r}^\alpha$ типа Лизоркина--Самко, 


связанные с комплексными\newline степенями телеграфного оператора}{\No\,7}


\tl{Ченцов А.Г.}


{К вопросу о построении корректных расширений


в классе конечно-\newline аддитивных мер}{\No\,2}


\tl{Чурбанов Ю.Д.}{Геометрия


однородных $\Phi$-пространств порядка 5}{\No\,5}


\tl{Широкова Е.А.}


{Получение классов данных для корректной постановки


обратной\newline краевой задачи путем перепараметризации}{\No\,4}


\tl{Шишкин Г.И.}


{Оптимальные по порядку скорости сходимости кусочно-равномерные


сетки\newline для сингулярно возмущенных уравнений конвекции--диффузии}{\No\,3}


\tl{Шпирко С.В.}


{О существовании и единственности решения


задачи равновесного\newline программирования}{\No\,12}


\tl{Щеглова А.А.}


{Линейные алгебро-дифференциальные системы


с переменным отклонением\newline аргумента}{\No\,6}


\tl{Щеглова А.А.}


{Об алгебро-дифференциальных системах с отклоняющимся аргументом}{\No\,7}


\tl{Эскин Л.Д.}  


{О некоторых решениях задачи о распаде разрыва в динамике  \newline


неньютоновской жидкости}{\No\,4}


\tl{Якубовский З.Ю., Старков В.В.}


{О совпадении двух линейно-инвариантных семейств\newline функций}{\No\,4}











\bigskip


{\centerline {КРАТКИЕ СООБЩЕНИЯ}}


\bigskip








\tl{Алшибая Э.Д.} 


{Об аффинных связностях на распределении 


гиперплоскостных элементов\newline в $A_{n+1}$}{\No\,8}


\tl{Афиногентова Е.В., Щенников В.Н.}


{Построение оценок погрешности линеаризации\newline систем 


конечно-разностных уравнений}{\No\,8}


\tl{Виноградова Г.Ю., Деундяк В.М.}


{Об индексе и гомотопической классификации \newline


семейств сингулярных операторов с


диэдральной группой сдвигов и $PC$-коэффициентами}{\No\,11}


\tl{Галимянов А.Ф.}


{К прямым методам решения интегральных уравнений 


с\newline логарифмически ослабленными ядрами Коши 


на разомкнутых контурах}{\No\,3}


\tl{Гарипов И.Б.} 


{Решение первой краевой задачи для $B$-параболического 


уравнения методом\newline потенциалов}{\No\,9}


\tl{Деундяк В.М., Виноградова Г.Ю.}


{Об индексе и гомотопической классификации \newline


семейств сингулярных операторов с


диэдральной группой сдвигов и $PC$-коэффициентами}{\No\,11}


\tl{Жогова Т.Б.} 


{Семейства $V^2_{2n-1}$, $U^m_{2n-1}$ и их проективное изгибание}{\No\,11}


\tl{Зинченко А.Б.} 


{Структура множества Парето некоторых векторных задач  


составления\newline расписаний}{\No\,7}


\tl{Кокурин М.Ю., Юсупова Н.А.} 


{О необходимых и достаточных условиях медленной\newline


сходимости методов решения линейных некорректных задач}{\No\,2}


\tl{Кравцов В.М., Кравцов М.К., Лукшин Е.В.}


{О типах $(3n-2)$-нецелочисленных\newline вершин многогранника трехиндексной


аксиальной задачи выбора}{\No\,12}


\tl{Кравцов М.К., Кравцов В.М., Лукшин Е.В.}


{О типах $(3n-2)$-нецелочисленных\newline вершин многогранника трехиндексной


аксиальной задачи выбора}{\No\,12}


\tl{Крепкогорский В.Л.}


{Интерполяция пространств Бесова.


Нормы, заданные с помощью\newline модуля непрерывностей}{\No\,1}


\tl{Лейнартас Д.Е.}


{О задаче Коши для многомерного разностного уравнения


с постоянными\newline коэффициентами}{\No\,1}


\tl{Лужецкая П.А., Ногин В.А.}


{$L_p\to L_q$-оценки для некоторых операторов типа потенциала\newline


с осциллирующими символами и их приложение}{\No\,11}


\tl{Лукшин Е.В., Кравцов В.М.m Кравцов М.К.}


{О типах $(3n-2)$-нецелочисленных\newline вершин многогранника трехиндексной


аксиальной задачи выбора}{\No\,12}


\tl{Мавлявиев Р.М.}


{Решение основных краевых задач для одного


$B$-эллиптического\newline уравнения методом потенциалов}{\No\,9}


\tl{Ногин В.А., Лужецкая П.А.}


{$L_p\to L_q$-оценки для некоторых операторов типа потенциала\newline


с осциллирующими символами и их приложение}{\No\,11}


\tl{Рахимов И.К.}


{Проекционные методы решения сингулярного интегрального уравнения \newline


Теодорсена}{\No\,9}


\tl{Романова С.В.} 


{Асимптотические оценки линейных функционалов для ограниченных \newline


функций, не принимающих нулевого значения}{\No\,11}


\tl{Скворцова Г.Ш.}


{О слабой секвенциальной полноте факторпространств


пространства\newline интегрируемых операторов}{\No\,9}


\tl{Хусаинова Э.Д.}


{О некоторых сингулярных задачах дифракции 


с общей границей}{\No\,9}


\tl{Цалюк З.Б., Чистоклетова С.С.}


{Асимптотика решений уравнений


Вольтерра--Гаммер-\newline штейна}{\No\,6}


\tl{Чистоклетова С.С., Цалюк З.Б.}


{Асимптотика решений уравнений


Вольтерра--Гаммер-\newline штейна}{\No\,6}


\tl{Шлапунов А.А.}


{О двойственности в пространствах полигармонических функций}{\No\,8}


\tl{Щенников В.Н., Афиногентова Е.В.}


{Построение оценок погрешности линеаризации\newline систем 


конечно-разностных уравнений}{\No\,8}


\tl{Эскин Л.Д.}


{Об интегральном уравнении, описывающем ориентационные


фазовые\newline переходы в модели Парсонса для системы осесимметричных частиц}{\No\,3}


\tl{Юсупова Н.А., Кокурин М.Ю.} 


{О необходимых и достаточных условиях медленной\newline


сходимости методов решения линейных некорректных задач}{\No\,2}





\medskip





\tl{Сосов Е.Н.}{Письмо в редакцию}{\No\,1}


\medskip











\noindent{}Рефераты статей, депонированных в ВИНИТИ


через журнал ``Известия вузов.\newline Математика''.\dotfill{\No\,4}








\medskip





\noindent{}Информация о научной конференции ``Геометрическая


теория функций и краевые\newline задачи'' \dotfill {\No\,10}








%% \tl{}{Математическая жизнь в России и зарубежом.}{\No\,11}











\end{document}








